\chapter{The Global Theorem}

\begin{lemma} \label{lem:hullscalarprod}
\leanok
\lean{GlobalTheorem.hull_scalar_prod}
Suppose $V = V_1, \ldots, V_m \subseteq \R^n$ be a finite sequence of points. Suppose $\hull V$ is its convex hull.
Let $S \in \hull V$ and $w \in \R^n$ be given. then
    \[
        \langle S ,w \rangle \leq \max_{i} \langle V_i ,w\rangle
    \]
\end{lemma}

\begin{proof}
\leanok
This is a mild generalization of \cite{polyhedron.without.rupert}, Lemma 18.

Since $S \in \hull V$, we have
\[ S = \sum_{j=1}^m \lambda_j V_j \]
for some $\lambda_1,\ldots,\lambda_m \in [0,1]$ with
\[ 1 = \sum_{j=1}^m \lambda_j  \]
Therefore
\[
\langle S ,w \rangle = \left\langle  \sum_{j=1}^m \lambda_j V_j ,w \right\rangle
=  \sum_{j=1}^m \lambda_j \left\langle   V_j ,w \right\rangle
\le   \sum_{j=1}^m \lambda_j \max_{i} \langle V_i ,w\rangle
\]
\[
=  \max_{i} \langle V_i ,w\rangle \sum_{j=1}^m \lambda_j
=  \max_{i} \langle V_i ,w\rangle
\]
as required.
\end{proof}


\begin{lemma} \label{lem:leq1}
\lean{GlobalTheorem.rotation_third_partials_bounded}
\leanok
    Let $S \in \R^3$ and $w \in \R^2$ be unit vectors and set $f(x_1,x_2,x_3) = \langle R(x_3) M(x_1,x_2)S,w \rangle$. Then for all $x_1,x_2,x_3 \in \R$ and any $i,j,k \in \{1,2,3\}$ it holds that
    \[
        \left|\frac{\dd^2 f}{\dd x_i \dd x_j}(x_1,x_2,x_3)\right|\leq 1
        \qquad\text{and}\qquad
        \left|\frac{\dd^3 f}{\dd x_i \dd x_j \dd x_k}(x_1,x_2,x_3)\right|\leq 1.
    \]
\end{lemma}

\begin{proof}\leanok
The second-partial bound is \cite{polyhedron.without.rupert}, Lemma 19.  The
third-partial bound follows by the same argument: every third partial of $f$
is $\langle A S, w\rangle$ for $A$ a composition of $\pm R$ or $\pm R'$ with a
matrix from the $M$-derivative family, all of operator norm at most one.
(Note $M^{\theta\theta\theta} = -M^\theta$ and $M^{\phi\phi\phi} = -M^\phi$,
so only the two mixed third derivatives $M^{\theta\theta\phi}$,
$M^{\theta\phi\phi}$ are genuinely new matrices.)
\end{proof}

\begin{lemma} \label{lem:n2}
\lean{GlobalTheorem.bounded_partials_control_difference2}
\leanok
    Let $f:\R^n\to \R$ be a $C^3$-function, let
    $\varepsilon_1,\dots,\varepsilon_n \geq 0$, and let
    $x_1,\dots,x_n,y_1,\dots,y_n \in \R$ be such that
    $|x_i-y_i|\leq \varepsilon_i$ for all $i$.
    If
    \(
        \left|\partial_{x_i}\partial_{x_j}\partial_{x_k}f(v)\right| \leq 1
    \)
    for all $i,j,k \in \{1,\dots,n\}$ and all $v \in \R^n$, then
    \[
    |f(x)-f(y)|\leq \sum_{i=1}^n \varepsilon_i |\partial_{x_i} f(x)|
      + \frac{1}{2} \sum_{i=1}^n \sum_{j=1}^n \varepsilon_i \varepsilon_j
          |\partial_{x_i}\partial_{x_j} f(x)|
      + \frac{1}{6}\Big(\sum_{i=1}^n \varepsilon_i\Big)^3.
    \]
    (At $\varepsilon_1 = \dots = \varepsilon_n = \varepsilon$ this recovers the
    isotropic remainder $\frac{n^3}{6}\varepsilon^3$.)
\end{lemma}

\begin{proof}
\leanok
This strengthens \cite{polyhedron.without.rupert}, Lemma 20, by one Taylor
order and per-axis radii: expand $g(t) = f((1-t)x + ty)$ to second order with
Lagrange remainder, bound $|g'(0)|$ by the first sum, $|g''(0)|/2$ by the
second (the second partials are evaluated exactly at $x$, not bounded), and
$|g'''(c)| \le (\sum_i \varepsilon_i)^3$ using the third-partial bound.
\end{proof}

\begin{lemma}
\label{lem:rotation_derivatives}
\leanok
\lean{GlobalTheorem.rotation_partials_exist, GlobalTheorem.rotation_partials_exist_outer,
 GlobalTheorem.partials_helper0, GlobalTheorem.partials_helper1, GlobalTheorem.partials_helper2,
 GlobalTheorem.partials_helper3, GlobalTheorem.partials_helper4}
The partial derivatives of all relevant rotations, projections, and inner products
used in the Global Theorem are as expected.
Specifically:
\begin{itemize}
\item \[  f^\alpha(\theta,\phi,\alpha) = \langle R'(\alpha) M(\theta, \phi) S, w \rangle \]
\item \[  f^\theta(\theta,\phi,\alpha) = \langle R(\alpha) M^\theta(\theta, \phi) S, w \rangle \]
\item \[  f^\phi(\theta,\phi,\alpha) = \langle R(\alpha) M^\phi(\theta, \phi) S, w \rangle \]
\item \[  g^\theta(\theta,\phi) = \langle  M^\theta(\theta, \phi) P, w \rangle \]
\item \[  g^\phi(\theta,\phi) = \langle  M^\phi(\theta, \phi) P, w \rangle \]
\end{itemize}
where $f(\theta,\phi,\alpha) =  \langle R(\alpha) M(\theta,\phi) S / \|S\|, w\rangle $
and $g(\theta,\phi) =  \langle M(\theta,\phi) P / \|P\|, w\rangle $.
\end{lemma}

\begin{proof}\leanok
By basic properties of derivatives.
\end{proof}

\begin{theorem}[Global Theorem] \label{thm:global}
\lean{GlobalTheorem.global_theorem}
\leanok
    Let $\PPP$ be a pointsymmetric convex polyhedron with radius $\rho =1$ and let $S \in \PPP$. Further let $\thetab_1,\phib_1,\thetab_2,\phib_2,\alphab \in \R$, let
    $\epsilon_\alpha, \epsilon_{\theta_1}, \epsilon_{\phi_1}, \epsilon_{\theta_2}, \epsilon_{\phi_2} \geq 0$
    be per-axis radii, and let $w\in\R^2$ be a unit vector. Denote $\Mib \coloneqq M(\thetab_1, \phib_1)$, $ \Miib \coloneqq M(\thetab_2, \phib_2)$ as well as $\Mib^{\theta} \coloneqq M^\theta(\thetab_1, \phib_1)$, $\Mib^{\phi} \coloneqq M^\phi(\thetab_1, \phib_1)$ and analogously for $\Miib^{\theta}, \Miib^{\phi}$. Finally set
    \begin{align*}
        G \coloneqq{}& \langle R(\alphab) \Mib S,w \rangle - \epsilon_\alpha|\langle R'(\alphab)  \Mib S,w \rangle| - \epsilon_{\theta_1}|\langle R(\alphab) \Mib^\theta S,w \rangle| - \epsilon_{\phi_1}|\langle R(\alphab) \Mib^\phi S,w \rangle|\\
        &- \frac{1}{2}\big(\epsilon_\alpha^2|\langle R(\alphab) \Mib S,w \rangle|
          + 2\epsilon_\alpha\epsilon_{\theta_1}|\langle R'(\alphab) \Mib^\theta S,w \rangle| + 2\epsilon_\alpha\epsilon_{\phi_1}|\langle R'(\alphab) \Mib^\phi S,w \rangle|\\
        &\qquad\quad + \epsilon_{\theta_1}^2|\langle R(\alphab) \Mib^{\theta\theta} S,w \rangle| + 2\epsilon_{\theta_1}\epsilon_{\phi_1}|\langle R(\alphab) \Mib^{\theta\phi} S,w \rangle| + \epsilon_{\phi_1}^2|\langle R(\alphab) \Mib^{\phi\phi} S,w \rangle|\big)\\
        &- \frac{(\epsilon_\alpha+\epsilon_{\theta_1}+\epsilon_{\phi_1})^3}{6},\\
        H_P \coloneqq{}& \langle \Miib P,w \rangle + \epsilon_{\theta_2}|\langle \Miib^\theta P,w \rangle|+\epsilon_{\phi_2}|\langle  \Miib^\varphi P,w \rangle|\\
        &+ \frac{1}{2}\big(\epsilon_{\theta_2}^2|\langle \Miib^{\theta\theta} P,w \rangle| + 2\epsilon_{\theta_2}\epsilon_{\phi_2}|\langle \Miib^{\theta\phi} P,w \rangle| + \epsilon_{\phi_2}^2|\langle \Miib^{\phi\phi} P,w \rangle|\big) + \frac{(\epsilon_{\theta_2}+\epsilon_{\phi_2})^3}{6}, \quad \text{ for } P \in \PPP.
    \end{align*}
    If $G>\max_{P\in \PPP} H_P$ then there does not exist a solution to Rupert's condition with
    $$(\theta_1,\varphi_1,\theta_2,\varphi_2,\alpha) \in U \coloneqq [\thetab_1\pm\epsilon_{\theta_1}]\times[\phib_1\pm\epsilon_{\phi_1}]\times[\thetab_2\pm\epsilon_{\theta_2}]\times[\phib_2\pm\epsilon_{\phi_2}]\times[\alphab\pm\epsilon_\alpha] \subseteq \R^5.$$

    (This is a second-order, anisotropic strengthening of
    \cite{polyhedron.without.rupert}, Theorem 17, whose penalty is first-order
    and isotropic with quadratic remainders $9\epsilon^2/2$ and $2\epsilon^2$:
    here the exact second partials at the center pose are charged with
    per-axis weights — with multiplicities from the symmetric Hessian table —
    and only the cubic Lagrange remainders
    $(\epsilon_\alpha+\epsilon_{\theta_1}+\epsilon_{\phi_1})^3/6$ and
    $(\epsilon_{\theta_2}+\epsilon_{\phi_2})^3/6$ are bounded via
    Lemma~\ref{lem:leq1}. On the diagonal
    $\epsilon_\alpha = \epsilon_{\theta_1} = \dots = \epsilon$, these are
    $\frac{3^3}{6}\epsilon^3 = \frac{9\epsilon^3}{2}$ and
    $\frac{2^3}{6}\epsilon^3 = \frac{4\epsilon^3}{3}$.)
\end{theorem}

\begin{proof}
\leanok
\uses{lem:hullscalarprod,lem:leq1,lem:n2,lem:rotation_derivatives}
See \cite{polyhedron.without.rupert}, Section~{4.2}.
\end{proof}
